\chapter{Assignment 2: worked solutions on functional calculus}

\begin{context}
After completing this assignment, you will be able to verify a concrete representation of $C_b(X)$, reconstruct
spectral characterizations of normal elements, and distinguish norm convergence from strong operator
convergence. The dependencies given after each solution show which hypotheses justify each step.
\end{context}

\section{Problem 2.1: \texorpdfstring{$C_b(X)$}{Cb(X)} as a concrete \texorpdfstring{$C^*$}{C-star}-algebra}

Let $X$ be locally compact and Hausdorff.


Let $H=\ell^2(X)$ and, for $f\in C_b(X)$, define
\[
 (M_f\xi)(x)=f(x)\xi(x).
\]
For an arbitrary index set, the sum in $\|\xi\|_2^2$ is understood as
\[
 \|\xi\|_2^2=\sup_{F\subset X\text{ finite}}\sum_{x\in F}|\xi(x)|^2.
\]
We establish the lower norm bound when $X\ne\varnothing$.

\noindent\textbf{Statement.} Prove that $M:C_b(X)\to\mathcal B(\ell^2(X))$, $f\mapsto M_f$, is an isometric $*$-homomorphism.

\begin{proof}[Solution]
If $X=\varnothing$, $C_b(X)$ contains only the empty function and $\ell^2(X)=\{0\}$; therefore
$\mathcal B(\ell^2(X))=\{0\}$ and the unique map $f\mapsto M_f$ preserves operations and
norms. Assume henceforth that $X\ne\varnothing$.

For $\xi\in\ell^2(X)$,
\[
\|M_f\xi\|_2^2
=\sum_{x\in X}|f(x)|^2|\xi(x)|^2
\le\|f\|_\infty^2\|\xi\|_2^2,
\]
so $\|M_f\|\le\|f\|_\infty$. For the reverse inequality, given
$\varepsilon>0$ we choose $x_0$ with $|f(x_0)|>\|f\|_\infty-\varepsilon$ and take the vector
$e_{x_0}$ supported at $x_0$. Then
$\|M_fe_{x_0}\|=|f(x_0)|$, whence $\|M_f\|=\|f\|_\infty$.

Additionally,
\[
M_{f+g}=M_f+M_g,
\qquad M_{fg}=M_fM_g,
\qquad M_1=I.
\]
Finally,
\[
\langle M_f\xi,\eta\rangle
=\sum_x f(x)\xi(x)\overline{\eta(x)}
=\langle \xi,M_{\bar f}\eta\rangle,
\]
so $M_f^*=M_{\bar f}$. Therefore $M$ is an isometric $*$-homomorphism.
\end{proof}

\proofdiagram
{$f\in C_b(X)$}
{Norm in $\ell^2(X)$; pointwise conjugation}
{$\|M_f\|=\|f\|_\infty$, $M_{fg}=M_fM_g$, $M_{\bar f}=M_f^*$}
{$f\mapsto M_f$ is an isometric $*$-homomorphism}

\begin{cthejemplo}{(A diagonal operator on the space of sequences)}
To see the representation of Problem 2.1 over a known space, let's take $X=\mathbb N$ and
$f(n)=1/(n+1)$.
In $H=\ell^2(\mathbb N)$ we have $M_fe_n=(n+1)^{-1}e_n$; also,
\[
 \|M_f\|=\sup_{n\ge1}\|M_fe_n\|=\frac12=\|f\|_\infty.
\]
The bounded function appears as a diagonal operator, and its norm is preserved by
the isometric $*$-homomorphism of Problem 2.1.
\end{cthejemplo}

\noindent\textbf{Dependencies.} The definition of $C_b(X)$ from Part 1, the norm on $\ell^2(X)$, and the Hilbert space characterization of the adjoint. Here $X$ serves as an index set; its topology is not used to define $\ell^2(X)$.

\noindent\textbf{How this problem is used later.} It exhibits an isometric $*$-monomorphism into $\mathcal B(H)$, a concrete model of the homomorphisms and representations studied in Section 2.

\section{Problem 2.2: spectra of idempotents, products, and unitaries}

Let $\A$ be a unital Banach algebra. For the conclusion about unitaries in part (c), $\A$
is also a $C^*$-algebra.

\subsection*{2.2(a) Spectrum of an idempotent}

\noindent\textbf{Statement.} If $e\in\A$ and $e^2=e$, prove $\spec(e)\subset\{0,1\}$.

\begin{proof}[Solution]
If $\lambda\neq0,1$, set
\[
 b=\frac1\lambda 1+\frac{1}{\lambda(\lambda-1)}e.
\]
Since $e(1-e)=(1-e)e=0$ and $e^2=e$, we can write
\[
 \lambda1-e=\lambda(1-e)+(\lambda-1)e,
 \qquad b=\lambda^{-1}(1-e)+(\lambda-1)^{-1}e.
\]
On multiplication, the cross terms vanish and
$(\lambda1-e)b=b(\lambda1-e)=(1-e)+e=1$. Thus every
$\lambda\notin\{0,1\}$ belongs to the resolvent set and
$\spec(e)\subset\{0,1\}$.
\end{proof}

\pagebreak
\proofdiagram
{$e^2=e$}
{Solve algebraically $(\lambda1-e)^{-1}$}
{For $\lambda\notin\{0,1\}$, construct $\lambda^{-1}1+[\lambda(\lambda-1)]^{-1}e$}
{$\spec(e)\subset\{0,1\}$}

\begin{cthejemplo}{(Spectrum of a matrix projection)}
In $M_2(\mathbb C)$ we want to identify the two points allowed by Problem 2.2(a).
For $P=\operatorname{diag}(1,0)$, $P^2=P$ and
$\lambda I-P=\operatorname{diag}(\lambda-1,\lambda)$.
This matrix is invertible exactly when $\lambda\ne0,1$; therefore
$\spec(P)=\{0,1\}$. If $P$ is replaced by $0$ or $I$, the same calculation leaves only
$\{0\}$ or $\{1\}$, respectively.
\end{cthejemplo}

\noindent\textbf{Dependencies.} The equation $e^2=e$ allows us to check an explicit inverse of $\lambda1-e$ for $\lambda\ne0,1$; the definition of spectrum in Part 1 applies.

\noindent\textbf{How this problem is used later.} This gives the forward implication of the spectral characterization of projections in Problem 2.6(c) and in Proposition 2.28.

\subsection*{2.2(b) Spectral radius of a commutative product}

\noindent\textbf{Statement.} If $x,y\in\A$ commute, prove $r(xy)\le r(x)r(y)$.

\begin{proof}[Solution]
Commutativity allows us to write $(xy)^n=x^ny^n$. So
\[
\|(xy)^n\|^{1/n}
\le\|x^n\|^{1/n}\|y^n\|^{1/n}.
\]
Passing to the limit and using the spectral radius formula,
$r(xy)\le r(x)r(y)$.
\end{proof}

\proofdiagram
{$xy=yx$}
{Spectral radius formula}
{$(xy)^n=x^ny^n$ and $\|(xy)^n\|\le\|x^n\|\|y^n\|$}
{$r(xy)\le r(x)r(y)$}

\begin{cthejemplo}{(The spectral radius inequality can be strict)}
In $C([0,1])$ we can check how much the bound of Problem 2.2(b) can be improved.
Let's take $f(t)=t$ and $g(t)=1-t$; they commute because multiplication is pointwise.
By Problem 1.10, $r(f)=\|f\|_\infty=1$ and $r(g)=1$, while
\[
 r(fg)=\max_{0\le t\le1}t(1-t)=\frac14<r(f)r(g)=1.
\]
The inequality is strict because the maxima of $f$ and $g$ occur at opposite endpoints.
\end{cthejemplo}

\noindent\textbf{Dependencies.} The commutation allows $(xy)^n=x^ny^n$ and the spectral radius formula from Part 1 converts the estimate of norms to the desired inequality.

\noindent\textbf{How this problem is used later.} It allows estimating products within commutative subalgebras such as $C^*(x,1)$. The classification of Problem 2.6 uses functional calculus and does not depend on this inequality.

\subsection*{2.2(c) Spectrum of an invertible with controlled norm}

\noindent\textbf{Statement.} If $z\in\GL(\A)$ and $\|z\|,\|z^{-1}\|\le1$, prove $\spec(z)\subset\mathbb T$; deduce the same inclusion for every unitary element of a unital $C^*$-algebra.

\begin{proof}[Solution]
If $\lambda\in\spec(z)$, then $|\lambda|\le r(z)\le\|z\|\le1$. Since $z$ is invertible,
$0\notin\spec(z)$. Additionally
\[
 \lambda\in\spec(z)\quad\Longleftrightarrow\quad
 \lambda^{-1}\in\spec(z^{-1}),
\]
because
$\lambda1-z=-\lambda z(\lambda^{-1}1-z^{-1})$. Therefore
$|\lambda|^{-1}\le r(z^{-1})\le1$, that is $|\lambda|\ge1$. Hence $|\lambda|=1$.

If $u$ is unitary in a $C^*$-algebra, $u^{-1}=u^*$ and
$\|u\|=\|u^*\|=1$, so the above result gives
$\spec(u)\subset\mathbb T$.
\end{proof}

\proofdiagrambranch
{$z\in\GL(\A)$ and $\|z\|\le1$}
{$\|z^{-1}\|\le1$}
{$\lambda\in\spec(z)\Rightarrow |\lambda|\le1$ and $\lambda^{-1}\in\spec(z^{-1})$}
{$|\lambda|^{-1}\le1$}
{$|\lambda|=1$, then $\spec(z)\subset\mathbb T$}

\begin{cthejemplo}{(Spectrum of a unitary diagonal matrix)}
To see the unitary conclusion of Problem 2.2(c) with two spectral values, let's take
$U=\operatorname{diag}(1,i)$ in $M_2(\mathbb C)$.
We have $U^*U=I$, $U^{-1}=U^*$ and $\|U\|=\|U^{-1}\|=1$.
Moreover, $\det(\lambda I-U)=(\lambda-1)(\lambda-i)$, so
$\spec(U)=\{1,i\}\subset\mathbb T$.
The norm bounds force both points to be on the unit circle.
\end{cthejemplo}

\noindent\textbf{Dependencies.} The bound $r(z)\le\|z\|$, the relationship between the spectra of an invertible and its inverse, and both hypotheses $\|z\|,\|z^{-1}\|\le1$ are used.

\noindent\textbf{How this problem is used later.} This gives the unitary implication of Problem 2.6(b) and Proposition 2.28; for a unitary $u$, the identity $C^*$ ensures $\|u\|=\|u^{-1}\|=1$.

\begin{investigacion}[title={Unitary rotations on $L^2(\mathbb T)$}]
To apply Problem 2.2(c) in the Hilbert space model of \cite{BFK2006}, fix $z\in\mathbb T$ and
define $(U_z\varphi)(t)=\varphi(zt)$ on $L^2(\mathbb T)$.
The invariance of the Haar measure under rotations gives $\|U_z\varphi\|_2=\|\varphi\|_2$;
also $U_{z^{-1}}U_z=U_zU_{z^{-1}}=I$. Thus $U_z$ is unitary and
$\|U_z\|=\|U_z^{-1}\|=1$.
By Problem 2.2(c), $\spec(U_z)\subset\mathbb T$. The conclusion uses those properties
verified on $L^2(\mathbb T)$.
\end{investigacion}

\pagebreak
\section{Problem 2.3: products of \texorpdfstring{$C^*$}{C-star}-algebras and the ideal \texorpdfstring{$c_0$}{c0}}

Let $I$ be a set of indices and $(A_i)_{i\in I}$ be a family of $C^*$-algebras. In part
(b), take $I=\mathbb N$.

\subsection*{2.3(a) Bounded product of a family of $C^*$-algebras}

\noindent\textbf{Statement.} Prove that the product of families $(a_i)_{i\in I}$ with $a_i\in A_i$ and $\sup_i\|a_i\|<\infty$ is a $C^*$-algebra with pointwise operations and supremum norm.

\begin{proof}[Solution]
Let
\[
\prod_{i\in I}^{\ell^\infty} A_i
=\{(a_i): a_i\in A_i,\ \sup_i\|a_i\|<\infty\}.
\]
Pointwise operations preserve uniform boundedness; in particular,
\[
 \sup_i\|a_i+b_i\|\le\sup_i\|a_i\|+\sup_i\|b_i\|,
 \qquad
 \sup_i\|a_ib_i\|\le\bigl(\sup_i\|a_i\|\bigr)
                        \bigl(\sup_i\|b_i\|\bigr).
\]
With the norm $\|(a_i)\|=\sup_i\|a_i\|$, this is a normed algebra. If
$(a^{(n)})$ is Cauchy, each coordinate $a_i^{(n)}$ converges in $A_i$ to an element $a_i$.
Fix $\varepsilon>0$. There is $N$ such that
$\sup_i\|a_i^{(n)}-a_i^{(m)}\|<\varepsilon$ for $n,m\ge N$. Letting $m\to\infty$
coordinate by coordinate, we get $\|a_i^{(n)}-a_i\|\le\varepsilon$ for all $i$ and $n\ge N$.
In particular $\sup_i\|a_i\|\le\|a^{(N)}\|+\varepsilon<\infty$, and the same estimate shows
$\|a^{(n)}-a\|\le\varepsilon$ for $n\ge N$. Therefore the bounded product is complete.
Defining $(a_i)^*=(a_i^*)$,
\[
\|(a_i)^*(a_i)\|
=\sup_i\|a_i^*a_i\|
=\sup_i\|a_i\|^2
=\|(a_i)\|^2.
\]
This is a $C^*$-algebra.
\end{proof}

\proofdiagram
{$\{A_i\}_{i\in I}$ $C^*$-algebras}
{Coordinatewise operations; supremum norm}
{Coordinate completeness and $\|(a_i)^*(a_i)\|=\sup_i\|a_i^*a_i\|$}
{The bounded product $\prod_i A_i$ is a $C^*$-algebra}

\begin{cthejemplo}{(A bounded product of matrices)}
In the product of Problem 2.3(a), each coordinate can be a different matrix.
Let's take $A_n=M_2(\mathbb C)$ and $a_n=\operatorname{diag}(1,1/n)$ for $n\ge1$.
The family $a=(a_n)$ belongs to the bounded product because $\sup_n\|a_n\|=1$; moreover,
\[
 a_n^*a_n=\operatorname{diag}(1,1/n^2),
 \qquad \|a^*a\|=1=\|a\|^2.
\]
The $C^*$-identity for the product is verified coordinate by coordinate, as in Problem 2.3(a).
\end{cthejemplo}

\noindent\textbf{Dependencies.} Coordinate-by-coordinate operations, completeness of each $A_i$, and the $C^*$-identity are applied on each coordinate. The supremum norm restricts the product to uniformly bounded families.

\noindent\textbf{How this problem is used later.} This constructs the ambient algebra of Problem 2.3(b), where families tending to zero form a closed ideal.

\subsection*{2.3(b) The sum $c_0$ is a closed ideal}

\noindent\textbf{Statement.} For $I=\mathbb N$, prove that the sequences $(a_n)$ of the bounded product satisfying $\|a_n\|\to0$ form a closed two-sided ideal.

\begin{proof}[Solution]
Let
\[
\bigoplus_{n\in\mathbb N}^{c_0}A_n
=\{(a_n)\in\prod^{\ell^\infty}A_n:\|a_n\|\to0\}.
\]
If $\|a_n\|,\|c_n\|\to0$, then
$\|a_n+c_n\|\le\|a_n\|+\|c_n\|\to0$ and
$\|\alpha a_n\|=|\alpha|\|a_n\|\to0$; therefore the $c_0$-sum is a linear subspace.
If $b=(b_n)$ is in the bounded product and $a=(a_n)$ is in the $c_0$-sum, then
\[
\|b_na_n\|\le\|b\|_\infty\|a_n\|\to0,
\quad
\|a_nb_n\|\le\|a_n\|\|b\|_\infty\to0,
\]
so it is a two-sided ideal. If $a^{(k)}\to a$ in supremum norm and each $a^{(k)}$ tends to zero in
coordinates, fix a positive $\varepsilon$ and choose $k$ with
$\|a-a^{(k)}\|_\infty<\varepsilon/2$ and then $N$ with
$\|a_n^{(k)}\|<\varepsilon/2$ for $n\ge N$. Then $\|a_n\|<\varepsilon$ for $n\ge N$.
Therefore the ideal is closed.
Furthermore $(a_n^*)$ belongs to the same ideal because $\|a_n^*\|=\|a_n\|\to0$.
\end{proof}

\proofdiagram
{$a=(a_n)$ with $\|a_n\|\to0$}
{Coordinatewise product and supremum norm}
{$\|b_na_n\|\le\|b\|_\infty\|a_n\|\to0$; uniform limit preserves condition}
{$\bigoplus_n A_n$ is a closed two-sided ideal}

\begin{cthejemplo}{(The classical ideal of null sequences)}
To recognize the construction of Problem 2.3(b), let's put $A_n=\mathbb C$ for each $n$.
The sequence $a_n=1/n$ belongs to $c_0$, while $b_n=(-1)^n$ is in $\ell^\infty$.
Their product satisfies $|b_na_n|=1/n\to0$, and the same goes for $a_nb_n$.
Thus the ideal $c_0\subset\ell^\infty$ is the scalar case of the two-sided ideal property proved in Problem 2.3(b).
\end{cthejemplo}

\noindent\textbf{Dependencies.} The bounded product of Problem 2.3(a) provides the supremum norm. Coordinate-by-coordinate estimates prove the ideal property, and the uniform criterion proves closedness.

\noindent\textbf{How this problem is used later.} This allows us to form the quotient of the product by the ideal $c_0$ using Problem 1.4; interpreting that quotient requires further results about quotients of $C^*$-algebras.

\section{Problem 2.4: continuous spectral mapping theorem}

Let $\A$ be a unital $C^*$-algebra, let $x\in\A$ be normal, and let $f:\spec(x)\to\mathbb C$ be continuous.

\noindent\textbf{Statement.} Prove the continuous spectral mapping theorem: $\spec_\A(f(x))=f(\spec_\A(x))$.

\begin{proof}[Solution]
Work in $B=C^*(x,1)$, which is commutative. By spectral permanence, the spectrum of
$f(x)$ is the same whether computed in $B$ or in $\A$. The functional calculus is a $*$-isomorphism
\[
 C(\spec(x))\longrightarrow B.
\]
Under this isomorphism, $f(x)$ corresponds to the function $f$. By Problem 1.10, the
spectrum of $f$ in $C(\spec(x))$ is its range. Therefore
\[
 \spec_\A(f(x))=\spec_B(f(x))=f(\spec(x)).
\]
\end{proof}

\proofdiagrambranch
{$x$ normal}
{$f\in C(\spec(x))$}
{Functional calculus + spectral permanence in $C^*(x,1)$}
{Identify $C^*(x,1)\cong C(\spec(x))$}
{$\spec(f(x))=f(\spec(x))$}

\begin{cthejemplo}{(Squaring a self-adjoint element using its spectrum)}
In $M_2(\mathbb C)$, take $x=\operatorname{diag}(-2,3)$ to check Problem 2.4.
The element is self-adjoint and $\spec(x)=\{-2,3\}$.
As $x^2=\operatorname{diag}(4,9)$,
\[
 \spec(x^2)=\{4,9\}=\{\lambda^2:\lambda\in\spec(x)\}.
\]
The square transforms each spectral value by $t\mapsto t^2$ and produces a nonnegative spectrum.
\end{cthejemplo}

\noindent\textbf{Dependencies.} The functional calculus for normal $x$ identifies $C^*(x,1)$ with $C(\spec(x))$; the spectral permanence of Section 2 and Problem 1.10 allow the spectrum of $f(x)$ to be calculated in the ambient algebra.

\noindent\textbf{How this problem is used later.} This transfers equalities and spectral restrictions from functions to operators; Problem 2.6 applies it to recognize self-adjoint elements, unitary elements, and normal projections.

\section{Problem 2.5: \texorpdfstring{$X\cong\Omega(C(X))$}{X is homeomorphic to Omega(C(X))}}

Let $X$ be a compact Hausdorff space and, for $x\in X$, let
$\operatorname{ev}_x(f)=f(x)$ be the evaluation map on $C(X)$.

\noindent\textbf{Statement.} Prove that $x\mapsto\operatorname{ev}_x$ defines a homeomorphism $X\to\Omega(C(X))$.

\begin{proof}[Solution]
For $f,g\in C(X)$, we have
$\operatorname{ev}_x(fg)=f(x)g(x)$ and $\operatorname{ev}_x(1)=1$;
also $|\operatorname{ev}_x(f)|\le\|f\|_\infty$. Thus each evaluation is a continuous character
and the map is well defined. It is injective because
if $x\neq y$, there is $f\in C(X)$ with $f(x)\neq f(y)$.

Now let $\varphi\in\Omega(C(X))$ and set $I=\ker\varphi$, a maximal ideal. We claim that there is
$x_0\in X$ such that $I\subset I_{x_0}:=\{f:f(x_0)=0\}$. Otherwise, for each $x\in X$ there would be
$f_x\in I$ with $f_x(x)\neq0$. The open sets
$U_x=\{y:|f_x(y)|>0\}$ would cover $X$; by compactness we could choose
$f_1,\dots,f_n\in I$ with $\sum_j|f_j|^2>0$ throughout $X$. Then
\[
 g:=\sum_{j=1}^n f_j^*f_j\in I
\]
would be a strictly positive function and therefore invertible in $C(X)$, impossible for a proper ideal. Thus there is $x_0$ with $I\subset I_{x_0}$. Since $I_{x_0}$ is proper, maximality gives
$I=I_{x_0}$. For each $f\in C(X)$, we have $f-f(x_0)1\in I=\ker\varphi$; since
$\varphi(1)=1$, it follows that $\varphi(f)=f(x_0)$. Therefore
$\varphi=\operatorname{ev}_{x_0}$.

The map $x\mapsto\operatorname{ev}_x$ is continuous for the weak-$*$ topology because, for each
$f\in C(X)$, the coordinate $x\mapsto\operatorname{ev}_x(f)=f(x)$ is continuous. As $X$ is compact
and $\Omega(C(X))$ is Hausdorff, a continuous bijection is a homeomorphism.
\end{proof}

\proofdiagrambranch
{$x\mapsto\operatorname{ev}_x$}
{Continuous functions separate points in compact Hausdorff spaces}
{For surjectivity, $\ker\varphi$ is a maximal ideal and must be $I_{x_0}$ for some $x_0$}
{Continuous bijection between compact Hausdorff spaces}
{$X\cong\Omega(C(X))$}

\begin{cthejemplo}{(A character given by a point on the circle)}
In $C(\mathbb T)$ we want to make the homeomorphism of Problem 2.5 visible.
Let's set $w=e^{i\theta}\in\mathbb T$ and consider the evaluation
$\operatorname{ev}_w(f)=f(w)$. For functions $h(z)=z$ and $g(z)=1+z$,
\[
 \operatorname{ev}_w(h)=w,\qquad \operatorname{ev}_w(g)=1+w.
\]
The value of the character on $h$ identifies the point $w$; Problem 2.5 ensures that every character in
$C(\mathbb T)$ arises from such an evaluation.
\end{cthejemplo}

\noindent\textbf{Dependencies.} Characters and maximal ideals from Part 1, point separation from $C(X)$, and compactness from $X$ are used. The common-zero-set argument uses $\sum_j|f_j|^2$.

\noindent\textbf{How this problem is used later.} This completes the uniqueness of the compact space in commutative Gelfand--Naimark: a $*$-isomorphism $C(X)\cong C(Y)$ induces $X\cong Y$.

\section{Problem 2.6: spectral characterizations and the role of normality}

Let $\A$ be a unital $C^*$-algebra and let $x\in\A$ be normal. Part (d) asks us to show that
normality cannot be omitted from the three converses.

\subsection*{2.6(a) Self-adjoint $\Longleftrightarrow$ real spectrum for normal elements}

\noindent\textbf{Statement.} Prove that $x=x^*$ if and only if $\spec(x)\subset\mathbb R$.

\begin{proof}[Solution]
The direct implication has already been proven in Section 2. For the converse, by functional calculus,
$x$ corresponds to the function $z$ and $x^*$ to $\overline z$ on $\spec(x)$. If the spectrum is real,
the functions agree, so $x=x^*$.
\end{proof}

\proofdiagram
{$x$ normal}
{Real spectrum of a self-adjoint element; functional calculus}
{If $\spec(x)\subset\mathbb R$, then $z=\overline z$ on the spectrum}
{$x=x^*$}

\begin{cthejemplo}{(Real spectrum without self-adjointness)}
The converse of Problem 2.6(a) requires normality; a triangular matrix allows you to see it.
For $N=\begin{pmatrix}0&1\\0&0\end{pmatrix}$,
$\det(\lambda I-N)=\lambda^2$, then $\spec(N)=\{0\}\subset\mathbb R$.
But $N^*=\begin{pmatrix}0&0\\1&0\end{pmatrix}\ne N$.
Additionally $N^*N=\operatorname{diag}(0,1)\ne\operatorname{diag}(1,0)=NN^*$.
A real spectrum does not force self-adjointness without normality, as shown in Problem 2.6(d).
\end{cthejemplo}

\noindent\textbf{Dependencies.} The reality of the spectrum of a self-adjoint element was proved in Section 2; the converse uses functional calculus in $C^*(x,1)$, available because $x$ is normal.

\noindent\textbf{How this problem is used later.} This identifies exactly which spectral condition characterizes self-adjointness. Problem 2.6(d) shows with a nilpotent matrix why the normality hypothesis is necessary.

\subsection*{2.6(b) Unitary $\Longleftrightarrow$ spectrum in $\mathbb T$ for normal elements}

\noindent\textbf{Statement.} Prove that $x$ is unitary if and only if $\spec(x)\subset\mathbb T$.

\begin{proof}[Solution]
If $x$ is unitary, Problem 2.2(c) gives $\spec(x)\subset\mathbb T$. Conversely, if $x$ is
normal and the spectrum lies in $\mathbb T$, then
\[
 x^*x=(\overline z z)(x)=(|z|^2)(x)=1.
\]
As $x$ is normal, $xx^*=x^*x=1$, so it is unitary.
\end{proof}

\proofdiagram
{$x$ normal}
{Problem 2.2(c); functional calculus}
{$\spec(x)\subset\mathbb T\Rightarrow |z|^2=1$ on spectrum}
{$x^*x=(\overline z z)(x)=1$ and, by normality, $xx^*=1$}

\begin{cthejemplo}{(Spectrum on the unit circle without unitarity)}
To show the limitation of the converse in Problem 2.6(b), let's take
$J=\begin{pmatrix}1&1\\0&1\end{pmatrix}$.
Its characteristic polynomial is $(\lambda-1)^2$, so
$\spec(J)=\{1\}\subset\mathbb T$, but
$J^*J=\begin{pmatrix}1&1\\1&2\end{pmatrix}\ne I$.
The spectrum is on the circle and the matrix is non-unitary: Problem 2.6(d) explains the role of normality.
\end{cthejemplo}

\noindent\textbf{Dependencies.} Problem 2.2(c) gives the forward implication. In the converse, normality allows functional calculus and the identity $x^*x=(|z|^2)(x)$.

\noindent\textbf{How this problem is used later.} It allows recognizing unitary elements from their spectrum, for example in algebras of operators with unitary shifts when their hypotheses have been verified.

\subsection*{2.6(c) Projection $\Longleftrightarrow$ spectrum on $\{0,1\}$ for normals}

\noindent\textbf{Statement.} Prove that $x$ is a projection if and only if $\spec(x)\subset\{0,1\}$.

\begin{proof}[Solution]
If $x$ is a projection, it is idempotent and Problem 2.2(a) gives
$\spec(x)\subset\{0,1\}$. For the converse, if $x$ is normal and its spectrum is contained in
$\{0,1\}$, functional calculus gives
\[
 x^2-x=(z^2-z)(x)=0,
 \qquad
 x^*-x=(\overline z-z)(x)=0.
\]
So $x$ is a projection.
\end{proof}

\proofdiagram
{$x$ normal}
{Spectrum of an idempotent; functional calculus}
{$z^2-z=0$ and $\overline z-z=0$ on $\{0,1\}$}
{$x^2=x$ and $x^*=x$}

\begin{cthejemplo}{(An idempotent with the expected spectrum that is not a projection)}
In Problem 2.6(c), self-adjointness is part of the definition of a projection. Let's consider
$E=\begin{pmatrix}1&1\\0&0\end{pmatrix}$.
A calculation gives $E^2=E$ and, by triangularity, $\spec(E)=\{0,1\}$; however,
$E^*=\begin{pmatrix}1&0\\1&0\end{pmatrix}\ne E$.
Thus, the permitted spectrum and idempotence are not enough without the normality used by the
converse of Problem 2.6(c).
\end{cthejemplo}

\noindent\textbf{Dependencies.} Problem 2.2(a) gives the forward implication. For the converse, normality allows $z^2=z$ and $\overline z=z$ to be transported from the spectrum to identities of $x$.

\noindent\textbf{How this problem is used later.} This connects the spectral classification of Section 2 with Cauchy projections only when, in addition to idempotence, the self-adjointness of those projections is verified.

\subsection*{2.6(d) Where is normality used?}

\noindent\textbf{Statement.} Display in $M_2(\mathbb C)$ three non-normal elements that invalidate the converses of parts (a)--(c) without normality, and locate the step of each proof that uses that hypothesis.

\begin{proof}[Solution]
Let's take
\[
 N=\begin{pmatrix}0&1\\0&0\end{pmatrix},
 \qquad J=I+N=\begin{pmatrix}1&1\\0&1\end{pmatrix},
 \qquad E=\begin{pmatrix}1&1\\0&0\end{pmatrix}.
\]
As they are triangular, $\spec(N)=\{0\}\subset\mathbb R$,
$\spec(J)=\{1\}\subset\mathbb T$ and $\spec(E)=\{0,1\}$.
However, $N\ne N^*$,
\[
 J^*J=\begin{pmatrix}1&1\\1&2\end{pmatrix}\ne I.
\]
$E$ is not a projection because, although $E^2=E$, it has $E^*\ne E$.
All three are non-normal: indeed,
$N^*N=\operatorname{diag}(0,1)\ne\operatorname{diag}(1,0)=NN^*$ and
$J^*J-JJ^*=N^*N-NN^*\ne0$; also,
$E^*E=\begin{pmatrix}1&1\\1&1\end{pmatrix}$ and
$EE^*=\operatorname{diag}(2,0)$ are different.

Normality was used in parts (a)--(c) so that $x$ and $x^*$ generate the
commutative $C^*$-algebra $C^*(x,1)$ and continuous functional calculus can be applied.
In (b) the equality $xx^*=x^*x$ was also used after obtaining $x^*x=1$.
\end{proof}

\proofdiagram
{$N$, $J=I+N$ and $E$ are non-normal}
{Continuous functional calculus was only built for normals}
{Functional calculus transfers equality of functions on the spectrum to equality of elements}
{Without normality, the spectrum does not determine these algebraic relationships}

\begin{cthejemplo}{(Two operators with the same spectrum)}
To see what the spectrum loses without normality, let's compare $I$ with
$J=\begin{pmatrix}1&1\\0&1\end{pmatrix}$ in $M_2(\mathbb C)$.
Both have spectrum $\{1\}$, but $I^*I=I$, while
$J^*J=\begin{pmatrix}1&1\\1&2\end{pmatrix}\ne I$.
Only $I$ is unitary. The concrete comparison explains why Problem 2.6(d) requires normality
before retrieving algebraic properties from the spectrum.
\end{cthejemplo}

\noindent\textbf{Dependencies.} Parts 2.6(a)--(c) use functional calculus on $C^*(x,1)$; that subalgebra is commutative because $x$ is normal. Jordan blocks provide counterexamples.

\noindent\textbf{How this problem is used later.} This delimits the scope of Proposition 2.28 and prevents us from inferring self-adjointness, unitarity or projection from the spectrum of a non-normal operator.

\pagebreak
\section{Problem 2.7: Unitary part of polar decomposition}

Let $\A$ be a unital $C^*$-algebra. We define $|a|$ for every $a\in\A$ and assume that $a$ is
invertible when constructing the unitary factor.


We define $|a|=(a^*a)^{1/2}$ by positive functional calculus.

\noindent\textbf{Statement.} For $a\in\GL(\A)$, prove that $u=a|a|^{-1}$ is unitary, where $|a|=(a^*a)^{1/2}$.

\begin{proof}[Solution]
If $a$ is invertible, then $a^*a$ has inverse $a^{-1}(a^*)^{-1}$.
Since $|a|^2=a^*a$ is invertible, so is $|a|$: the elements $|a|$ and
$(|a|^2)^{-1}$ commute, and $|a|(|a|^2)^{-1}$ is a two-sided inverse of $|a|$.
We can, therefore, define $u=a|a|^{-1}$.
As $|a|$ is positive, $|a|^*=|a|$ and
\[
 u^*=|a|^{-1}a^*.
\]
Then
\[
 u^*u=|a|^{-1}a^*a|a|^{-1}
 =|a|^{-1}|a|^2|a|^{-1}=1.
\]
Moreover,
\[
 uu^*=a|a|^{-2}a^*
 =a(a^*a)^{-1}a^*
 =a\,a^{-1}(a^*)^{-1}a^*=1.
\]
Thus $u$ is unitary and $a=u|a|$.
\end{proof}

\proofdiagram
{$a\in\GL(\A)$}
{Positive root $|a|=(a^*a)^{1/2}$}
{Define $u=a|a|^{-1}$ and calculate $u^*u$, $uu^*$}
{$u^*u=uu^*=1$; $u$ is unitary}

\begin{cthejemplo}{(Polar decomposition of a diagonal matrix)}
To apply the formula in Problem 2.7, let's take the invertible matrix
$A=\operatorname{diag}(2i,3)$.
Then $A^*A=\operatorname{diag}(4,9)$,
$|A|=\operatorname{diag}(2,3)$ and
\[
 U=A|A|^{-1}=\operatorname{diag}(i,1),\qquad U^*U=UU^*=I.
\]
The equality $A=U|A|$ separates a unitary part from a positive part, as stated in Problem 2.7.
\end{cthejemplo}

\noindent\textbf{Dependencies.} The positive root $|a|=(a^*a)^{1/2}$ comes from the functional calculus of Section 2. Invertibility of $a$ implies invertibility of both $a^*a$ and $|a|$.

\noindent\textbf{How this problem is used later.} This constructs the polar decomposition $a=u|a|$ of an invertible element and clarifies where invertibility is used to obtain a unitary factor in the same algebra.

\Needspace{15\baselineskip}
\section{Problem 2.8: continuity of functional calculus under convergence of normal elements}

Let $\A$ be a unital $C^*$-algebra and let $(x_n)$ be a sequence of normal elements
converging in norm to $x$. Fix a compact set $K\subset\mathbb C$ containing $\spec(x_n)$
for every $n$ and $\spec(x)$, and let $f\in C(K)$.

\noindent\textbf{Statement.} Prove that $x$ is normal and that $f(x_n)\to f(x)$ in norm for all $f\in C(K)$.

\begin{proof}[Solution]
By Problem 1.7, $x$ is normal. Let $f\in C(K)$ and $\varepsilon>0$. Stone--Weierstrass provides
a polynomial $p(z,\overline z)$ such that
\[
 \sup_{z\in K}|f(z)-p(z,\overline z)|<\varepsilon.
\]
Functional calculus is contractive, so for all $n$,
\[
\|f(x_n)-p(x_n,x_n^*)\|<\varepsilon,
\qquad
\|f(x)-p(x,x^*)\|<\varepsilon.
\]
Since $x_n\to x$, we also have $x_n^*\to x^*$; continuity of addition and multiplication gives
$p(x_n,x_n^*)\to p(x,x^*)$. Therefore, for large $n$,
\[
\|f(x_n)-f(x)\|
\le 2\varepsilon+\|p(x_n,x_n^*)-p(x,x^*)\|<3\varepsilon.
\]
As $\varepsilon$ is arbitrary, $f(x_n)\to f(x)$.
\end{proof}

\proofdiagrambranch
{$x_n\to x$ and each $x_n$ is normal}
{$\spec(x_n),\spec(x)\subset K$ compact}
{Closedness of normal elements; Stone--Weierstrass approximates $f$ by polynomials in $z,\bar z$}
{Polynomials are continuous in $x_n,x_n^*$; control error uniformly}
{$f(x_n)\to f(x)$}

\begin{cthejemplo}{(Continuity of the positive square root)}
To see Problem 2.8 in positive matrices, let's set
$x_n=\operatorname{diag}(1+1/n,4)$ and $x=\operatorname{diag}(1,4)$.
All these spectra lie in $K=[1,5]$ and $x_n\to x$ in norm.
With $f(t)=\sqrt t$,
\[
 f(x_n)=\operatorname{diag}(\sqrt{1+1/n},2)
 \longrightarrow\operatorname{diag}(1,2)=f(x).
\]
The calculation shows that the positive square root varies continuously under the assumptions of Problem 2.8.
\end{cthejemplo}

\noindent\textbf{Dependencies.} Problem 1.7 guarantees that $x$ is normal; Stone--Weierstrass approximates $f$ uniformly by polynomials in $z,\overline z$, and functional calculus controls the error.

\noindent\textbf{How this problem is used later.} Shows the stability in norm of $f(x_n)$, in particular of positive roots when the spectra are in one compact subset of $[0,\infty)$.

\Needspace{17\baselineskip}
\section{Problem 2.9: monotone convergence of self-adjoint operators}

Let $H$ be a Hilbert space. A net is indexed by a directed set $\Lambda$.
For self-adjoint $R,U\in\mathcal B(H)$, we write $R\le U$ when $U-R$ is positive,
that is, when $\langle(U-R)\xi,\xi\rangle\ge0$ for all $\xi\in H$.


Let $(T_\lambda)_{\lambda\in\Lambda}$ be a bounded net of self-adjoint operators with
$T_{\lambda_1}\le T_{\lambda_2}$ if $\lambda_1\le\lambda_2$.
Here ``pointwise convergence'' means that, for each $\xi\in H$,
$\|T_\lambda\xi-T\xi\|\to0$; it is strong operator convergence.

\noindent\textbf{Statement.} Prove that there exists a self-adjoint $T\in\mathcal B(H)$ such that $\|T_\lambda\xi-T\xi\|\to0$ for each $\xi\in H$.

\begin{proof}[Solution]
Let $M=\sup_\lambda\|T_\lambda\|$. For each $\xi\in H$, the numbers
$q_\lambda(\xi)=\langle T_\lambda\xi,\xi\rangle$ form an increasing net bounded by
$M\|\xi\|^2$, so they converge to a limit $q(\xi)$. Using the polarization identity we obtain
a bounded sesquilinear form
\[
 B(\xi,\eta)=\lim_\lambda\langle T_\lambda\xi,\eta\rangle.
\]
Indeed, $|\langle T_\lambda\xi,\eta\rangle|\le M\|\xi\|\|\eta\|$, and the same bound passes to the
limit.
With the convention that the inner product is linear in the first variable, the polarization identity
expresses this limit using four values of $q$:
\[
 B(\xi,\eta)=\frac14\bigl(q(\xi+\eta)-q(\xi-\eta)
   +i q(\xi+i\eta)-i q(\xi-i\eta)\bigr).
\]
By the Riesz representation theorem there is a unique $T\in\mathcal B(H)$ with
$B(\xi,\eta)=\langle T\xi,\eta\rangle$, and the bound of $B$ implies $\|T\|\le M$.
Since each $T_\lambda$ is self-adjoint, so is $T$, and
$T_\lambda\le T$.

Set $A_\lambda=T-T_\lambda\ge0$. Then
$\|A_\lambda\|\le 2M$. Since
$\spec(A_\lambda)\subset[0,\|A_\lambda\|]$, the function
$t\mapsto t(\|A_\lambda\|-t)$ is nonnegative on that spectrum. Positive functional calculus
gives $0\le A_\lambda^2\le\|A_\lambda\|A_\lambda\le2M A_\lambda$. Therefore
\[
\|(T-T_\lambda)\xi\|^2
=\langle A_\lambda^2\xi,\xi\rangle
\le2M\langle A_\lambda\xi,\xi\rangle.
\]
The last term tends to zero by definition of $T$. So
$T_\lambda\xi\to T\xi$ in norm for each $\xi$.
\end{proof}

\proofdiagrambranch
{Monotonicity: increasing $\langle T_\lambda\xi,\xi\rangle$}
{Uniform boundedness $\sup\|T_\lambda\|<\infty$}
{Define a limit form and use Riesz to obtain $T$}
{Use $0\le(T-T_\lambda)^2\le\|T-T_\lambda\|(T-T_\lambda)$}
{$T_\lambda\xi\to T\xi$ for each $\xi$}

\begin{cthejemplo}{(Strong convergence of increasing matrices)}
To see the limit in Problem 2.9, take in $\mathbb C^2$
$T_n=\operatorname{diag}(1-1/n,2-1/n)$, $n\ge1$.
We have $T_{n+1}-T_n=(1/n-1/(n+1))I\ge0$ and $\|T_n\|\le2$.
With $T=\operatorname{diag}(1,2)$,
\[
 \|(T-T_n)\xi\|=\frac1n\|\xi\|\longrightarrow0
 \qquad(\xi\in\mathbb C^2).
\]
The convergence here is even in operator norm; Problem 2.9 ensures strong convergence
for bounded nets on any Hilbert space.
\end{cthejemplo}

\noindent\textbf{Dependencies.} Convergence of bounded monotone nets of real numbers, polarization identity, Riesz representation, and positive functional calculus are used to justify $0\le A^2\le\|A\|A$ when $A\ge0$.

\noindent\textbf{How this problem is used later.} This gives strong convergence: $\|T_\lambda\xi-T\xi\|\to0$ for each $\xi\in H$. This result prepares techniques for operator topologies, but is not a consequence of the convergence in norm of Problem 2.8.

\section{Problem 2.10: spectrum of the unilateral shift}


Let $S:\ell^2(\mathbb N)\to\ell^2(\mathbb N)$ be given by
\[
S(x_1,x_2,\dots)=(0,x_1,x_2,\dots).
\]

\noindent\textbf{Statement.} Prove that the spectrum of the unilateral shift is $\spec(S)=\overline{\mathbb D}$.

\begin{proof}[Solution]
Since $S$ is an isometry, $\|S\|=1$ and by the spectral bound
$\spec(S)\subset\overline{\mathbb D}$. Now let $|\lambda|<1$. The vector
\[
 v_\lambda=(1,\overline\lambda,\overline\lambda^2,\ldots)
\]
belongs to $\ell^2$ and satisfies
\[
 S^*v_\lambda=\overline\lambda\,v_\lambda.
\]
Thus $\overline\lambda I-S^*$ is not invertible. If $\lambda I-S$ were invertible, then its
adjoint $\overline\lambda I-S^*$ would also be invertible, contradiction. Therefore
$\lambda\in\spec(S)$ for all $|\lambda|<1$. Since the spectrum is closed,
$\overline{\mathbb D}\subset\spec(S)$. Along with the opposite inclusion,
\[
 \spec(S)=\overline{\mathbb D}.
\]
\end{proof}

\proofdiagrambranch
{$\|S\|=1\Rightarrow\spec(S)\subset\overline{\mathbb D}$}
{For $|\lambda|<1$, $S^*$ has eigenvector of value $\overline\lambda$}
{If $\lambda I-S$ were invertible, its adjoint would be too}
{$\mathbb D\subset\spec(S)$; the spectrum is closed}
{$\spec(S)=\overline{\mathbb D}$}

\begin{cthejemplo}{(A spectral disk with no eigenvalues)}
The unilateral shift of Problem 2.10 illustrates the difference between spectrum and eigenvalues.
If $S\xi=\lambda\xi$, the first coordinate gives $0=\lambda\xi_1$; for
$\lambda\ne0$, the subsequent coordinates successively force $\xi_n=0$.
For $\lambda=0$, injectivity of $S$ likewise gives $\xi=0$.
Thus $S$ has no eigenvalues, although Problem 2.10 proves
$\spec(S)=\overline{\mathbb D}$.
\end{cthejemplo}

\noindent\textbf{Dependencies.} The spectral bound of Part 1 gives $\spec(S)\subset\overline{\mathbb D}$. Problem 1.11(d) already shows $0\in\spec(S)$; the eigenvectors of $S^*$ give all the interior points and the closure of the spectrum adds the boundary.

\noindent\textbf{How this problem is used later.} This contrasts the spectrum of a non-normal operator with the characterization of unitary operators in Problem 2.6(b). It does not allow attributing the same spectral disk to the composition shifts of the Haseman problem.
